Du point où le soleil se lève\\
Jusqu’aux limites de la terre,\\
Chantons le Christ notre prince,\\
Né de la Vierge Marie.\\
\\
Le bienheureux créateur du monde\\
Revêt un corps d’esclave ;\\
Par sa chair il libère toute chair\\
Afin de ne pas perdre sa créature.\\
\\
La grâce du ciel pénètre\\
Le sein maternel scellé ;\\
Le ventre d’une vierge porte des mystères\\
Qu’elle ne connaissait pas.\\
\\
La demeure de son coeur très pur\\
Devient soudain le temple de Dieu ;\\
Sans le contact d’aucun homme,\\
D’une parole elle conçoit son Fils.\\
\\
La Mère met au monde\\
Celui que Gabriel avait annoncé,\\
Que, sautant dans le ventre de sa mère,\\
Le Baptiste reconnaissait de son enclos.\\
\\
Il a supporté de coucher sur la paille,\\
Il n’a pas refusé la crèche ;\\
Il s’est nourri d’un humble lait,\\
Lui qui rassasie même les oiseaux.\\
\\
Les choeurs d’en-haut se réjouissent,\\
Et les anges chantent Dieu ;\\
Le pasteur, créateur de tout,\\
Se montre à des pasteurs.\\
\\
Gloire à toi, Jésus,\\
Qui es né de la Vierge,\\
Comme au Père et à l’Esprit bienfaisant\\
Dans les siècles éternels.